%=============================================================================
\section{实用技巧}
%=============================================================================

\subsection{利用 preempt 分区跑更多实验}

\texttt{preempt} 分区的资源上限更高（20 GPU vs 10 GPU），但作业可能被节点所有者的任务\textbf{抢占}——SLURM 会在约 30 秒内 kill 你的进程。

应对策略是\textbf{让你的程序支持断点续训}：

\begin{lstlisting}[language=python,title={在 train.py 中添加 checkpoint 逻辑}]
import os, torch

ckpt_path = f"{args.output_dir}/checkpoint.pt"

# 启动时检查是否有已保存的进度
if os.path.exists(ckpt_path):
    checkpoint = torch.load(ckpt_path)
    model.load_state_dict(checkpoint["model"])
    start_epoch = checkpoint["epoch"]
    print(f"Resumed from epoch {start_epoch}")
else:
    start_epoch = 0

# 训练过程中定期保存
for epoch in range(start_epoch, total_epochs):
    train_one_epoch(model, ...)
    if epoch % 10 == 0:  # 每 10 个 epoch 保存一次
        torch.save({"model": model.state_dict(), "epoch": epoch}, ckpt_path)
\end{lstlisting}

这样即使作业被抢占，重新提交后也能从上次的进度继续，而不是从头开始。

\subsection{用邮件通知监控实验}

\begin{lstlisting}[language=slurm]
#SBATCH --mail-type=END,FAIL
#SBATCH --mail-user=your.email@tufts.edu
\end{lstlisting}

对于 Array Job（一次提交上百个任务），建议只在失败时通知，否则你的邮箱会被淹没：

\begin{lstlisting}[language=slurm]
#SBATCH --mail-type=FAIL
\end{lstlisting}

\subsection{汇总实验结果}

所有实验跑完后，用一个简单的 Python 脚本汇总各组的结果：

\begin{lstlisting}[language=python,title={\texttt{collect\_results.py}}]
import os, json, pandas as pd

rows = []
for env in open("configs/envs.txt"):
    env = env.strip()
    for arch in open("configs/archs.txt"):
        arch = arch.strip()
        for seed in range(5):
            result_file = f"results/{env}/{arch}/seed_{seed}/metrics.json"
            if os.path.exists(result_file):
                with open(result_file) as f:
                    metrics = json.load(f)
                rows.append({"env": env, "arch": arch, "seed": seed,
                             **metrics})

df = pd.DataFrame(rows)
summary = df.groupby(["env", "arch"]).agg(
    reward_mean=("final_reward", "mean"),
    reward_std=("final_reward", "std"),
).reset_index()
summary.to_csv("results/summary.csv", index=False)
print(summary.to_string())
\end{lstlisting}

\subsection{检查哪些实验失败了}

\begin{lstlisting}[language=bash]
# 查看有错误输出的日志
grep -l "Error\|Traceback\|CANCELLED" logs/*.err

# 查看哪些结果文件缺失
python -c "
import os
envs = open('configs/envs.txt').read().split()
archs = open('configs/archs.txt').read().split()
for e in envs:
    for a in archs:
        for s in range(5):
            p = f'results/{e}/{a}/seed_{s}/metrics.json'
            if not os.path.exists(p):
                print(f'MISSING: {e} / {a} / seed {s}')
"
\end{lstlisting}
